% Copyright (c) 2026 Geovana Neves. Licensed under CC BY 4.0 (https://creativecommons.org/licenses/by/4.0/). See LICENSE-AND-AUTHORSHIP.md.
%
% workspaces/text/02-plan-run-follow-collect.tex
%
% Parts 3 to 6 of guide 2: plan, run, follow, collect and post.

\section{Plan}

\begin{frame}[fragile]{What plan does}
\begin{lstlisting}
cfd-matrix plan polar
\end{lstlisting}
\begin{itemize}
  \item \code{polar} is \code{inputs/matrices/polar.fs}; the \code{.fs} is
        optional.
  \item Plan reads the matrix and every file it names, resolves the air of
        each row, and asks the solver's adapter to check the inputs.
  \item It lists each point on one line and ends with the receipt and the
        next command; \code{--verbose} adds the files each run would write.
  \item It \textbf{launches nothing} and creates no point folder.
  \item It saves the plan receipt, \code{post/polar/plan.json}: the matrix
        and every input, as planned.
\end{itemize}
\begin{takeaway}Plan again after every edit of an input. \code{run} refuses a
matrix whose inputs changed since its receipt was written.\end{takeaway}
\end{frame}

\begin{frame}[fragile]{Planning part of a matrix}
\begin{lstlisting}
cfd-matrix plan polar --sims 1001 1002
cfd-matrix plan polar --sims 1001,1002
cfd-matrix plan polar --points 1001/DP-M850RE500AL+025
\end{lstlisting}
\begin{itemize}
  \item \code{--sims} takes \code{POL} values, separated by spaces or commas.
  \item \code{--points} takes point names, with or without the \code{POL/}
        prefix.
  \item A name the matrix does not carry is refused before anything
        happens.
  \item \code{run}, \code{collect} and \code{post} take the same two
        options.
\end{itemize}
\begin{block}{How much a command prints}
\code{--log-level quiet}, \code{normal} (the default) or \code{verbose};
\code{--verbose} is the short form of the last. \code{--warning-off}, written
before the command name, keeps warnings out of the console and in
\code{logs/warnings.log}.
\end{block}
\end{frame}

\section{Run}

\begin{frame}[fragile]{What run does}
\begin{lstlisting}
cfd-matrix run polar
\end{lstlisting}
\begin{itemize}
  \item Runs every active point of the planned matrix: on this machine, or
        as cluster jobs when the workspace has an HPC profile (Part 7).
  \item For each point it writes the folder
        \code{sims/sim_<POL>/datapoints/DP-<point name>/} with the solver's
        input deck, the link to the mesh, the solver's own output,
        \code{stdout.log} and \code{stderr.log}.
  \item A run on this machine also leaves \code{execution.json}, the record
        of the process that ran.
  \item It records every point in \code{runs.json}: inputs and hashes,
        status, start and finish, verdict.
\end{itemize}
\begin{cfdblimit}[The optional host guard]
\small The host guard is off by default. Add \code{--one-run-at-a-time} to
exclude other guarded local runs. Executable and memory checks always apply.
\end{cfdblimit}
\end{frame}

\begin{frame}{Running a matrix again}
\relax{\cfdbbodysize
\begin{tabularx}{\linewidth}{@{}P{0.33\linewidth}R@{}}
\toprule
\textbf{You type} & \textbf{What happens to a point already recorded} \\
\midrule
\code{run polar} & it refuses the run: finished evidence is not replaced by
  accident \\
\code{run polar --resume} & it is skipped, a failed one with the
  \code{--force-rerun} that redoes it, and the points not yet recorded run.
  A matrix that grew a row runs only the new row \\
\code{run polar --force-rerun POINT} & that point is archived, then run
  again. Repeat the option for several points \\
\code{run polar --force-rerun-all} & every recorded point of the matrix is
  archived, then run again \\
\bottomrule
\end{tabularx}}
\vspace{2mm}
\begin{itemize}\small
  \item Archived means: \code{runs.json} is copied to
        \code{archive/<stamp>/}, the point folder moves there, and the old
        record joins the history of the point. Nothing is deleted.
  \item A submitted cluster job is not canceled by a rerun: cancel it in
        the scheduler first.
\end{itemize}
\end{frame}

\section{Follow}

\begin{frame}{Six commands that only read}
\relax{\cfdbbodysize
\begin{tabularx}{\linewidth}{@{}P{0.34\linewidth}R@{}}
\toprule
\textbf{Command} & \textbf{Prints} \\
\midrule
\code{cfd-matrix status} & each simulation with the states of its points.
  \code{--points}: one row per point. \code{--failed}: only the simulations
  with a failed point \\
\code{cfd-workspace show TARGET} & one simulation, point or run record in full \\
\code{cfd-workspace history TARGET} & the current and the earlier records of a
  simulation or point, with their archives \\
\code{cfd-workspace log} & the recorded events in time order.
  \code{--problems}: the failures. \code{--since 2h}: the recent ones \\
\code{cfd-workspace trace PRODUCT} & the runs, inputs and hashes behind a
  product file \\
\code{cfd-workspace diff RUN_A RUN_B} & two run records, field by field \\
\bottomrule
\end{tabularx}}
\vspace{2mm}
\begin{itemize}\small
  \item All six take \code{--json} or \code{--csv}, to hand the rows to a
        program.
  \item They read the records in the workspace. None of them asks a
        scheduler or touches a solver.
\end{itemize}
\end{frame}

\begin{frame}{Status and verdict are two things}
\begin{columns}[T]
\begin{column}{0.48\textwidth}
\begin{block}{Status: where the point is}
\begin{itemize}\small
  \item \code{PLANNED}: registered, not started.
  \item \code{SUBMITTED}: handed to a scheduler.
  \item \code{RUNNING}: started on this machine.
  \item \code{COMPLETED}: the run ended and was read.
\end{itemize}
\end{block}
\end{column}
\begin{column}{0.48\textwidth}
\begin{block}{Verdict: how it ended}
\begin{itemize}\small
  \item \code{converged}
  \item \code{iteration_limit}
  \item \code{failed}
\end{itemize}
\end{block}
\end{column}
\end{columns}
\begin{cfdbpitfall}[Completed is not converged]
A completed point may have failed: the console shows it as
\code{FAILED_EXECUTION} or another \code{FAILED_} word, the records and tables
as \code{COMPLETED} with the verdict \code{failed}. Read the verdict, which
every table and every \code{status --points} row carries, before using a number.
\end{cfdbpitfall}
\end{frame}

\section{Collect and post}

\begin{frame}[fragile]{What collect writes}
\begin{lstlisting}
cfd-matrix collect polar
\end{lstlisting}
\relax{\cfdbtablesize
\begin{tabularx}{\linewidth}{@{}P{0.42\linewidth}R@{}}
\toprule
\textbf{File under} \code{post/polar/} & \textbf{Holds} \\
\midrule
\code{campaign_summary.json} & every recorded point with its result \\
\code{polars/FAM-ACOEF-P<POL>.csv} & coefficients per point and family,
  in wind, stability, body and rotor frames \\
\code{polars/FAM-DIM-P<POL>.csv} & dimensional loads per point and family \\
\code{polars/P<POL>-TOTAL-ACOEF.csv} & total coefficients per point \\
\code{polars/P<POL>.txt} & a 2D SU2 or OpenFOAM ALPHA sweep, XFOIL layout \\
\code{polars/P<POL>-<alias>_rotor.csv} & the rotor table, one row per point
  that turns the rotor \\
\code{polars/P<POL>_<matrix>_uns_avg.csv} & the unsteady polar: the mean over
  the averaging window of each time-accurate point \\
\code{sections/}, \code{fields/}, \code{probes/} & the products the post file
  declares, each in a folder \code{P<POL>/} \\
\code{products.json}, \code{post.log} & every product written, and every one
  skipped with the reason \\
\bottomrule
\end{tabularx}}
\begin{itemize}\small
  \item A value a point does not have is \code{NA}. No cell is quoted and
        none holds a comma, so any CSV reader takes the table.
  \item Without a matrix name, \code{collect} works on every recorded
        matrix.
  \item \code{--archive} keeps a copy of each replaced file under
        \code{archive/<stamp>/}.
\end{itemize}
\end{frame}

\begin{frame}{Products and boundary force sets}
\begin{itemize}\small
  \item Every polar has family files and family-based group files for
        \code{TOTAL}, \code{AIRFRAME} and each rotor that has members.
        Rotor solvers also write \code{FAM-RCOEF}; sectional distributions
        go under \code{sections/P<POL>/}, a volume cut under
        \code{fields/P<POL>/} and probes under \code{probes/P<POL>/}.
  \item The thrust and torque of the rotor table are counted along the
        declared shaft, so the sign of \code{CT} follows the axis. Where the
        solver supplies wall shear, the sectional loads are totals, pressure
        plus friction.
  \item Tables start with the canonical run block: conditions, air state,
        references, rotor data and averaging window. A value that does not
        apply is \code{NA}. Verdicts and histories live in the campaign summary.
  \item A setup \code{[groups]} table selects SU2 or OpenFOAM boundary markers.
        Its \code{force_groups}, including \code{TOTAL}, stay in the summary;
        they do not add boundary-group CSV files. Groups may overlap.
  \item \code{products.json}, schema version 2, states kinds, units, frames,
        families, groups, skips and SHA-256 digests. Edited or unrecorded
        products are protected against replacement. Preserve older products
        in a separate folder before recollecting.
\end{itemize}
\end{frame}

\begin{frame}{Collect or post}
\begin{columns}[T]
\begin{column}{0.48\textwidth}
\begin{block}{\code{cfd-matrix collect}}
\begin{itemize}\small
  \item Reads the saved results of each point.
  \item Picks up cluster jobs that finished.
  \item Updates the records, then rebuilds the tables.
\end{itemize}
\end{block}
\end{column}
\begin{column}{0.48\textwidth}
\begin{block}{\code{cfd-matrix post}}
\begin{itemize}\small
  \item Rebuilds the same tables from the records and the saved outputs
        alone.
  \item Runs nothing and asks nothing of a cluster.
\end{itemize}
\end{block}
\end{column}
\end{columns}
\vspace{1mm}
\relax{\cfdbbodysize
\begin{tabularx}{\linewidth}{@{}P{0.16\linewidth}R@{}}
\toprule
\textbf{Exit code} & \textbf{Meaning} \\
\midrule
0 & done \\
1 & done, and at least one point failed (\code{diff}: the two runs differ) \\
2 & refused: the message names the file and the cell or key at fault, and
  nothing was written \\
3 & done, and submitted work is still outstanding \\
\bottomrule
\end{tabularx}}
\end{frame}
